% TOP_PROBLEM: {"id": 3000037, "title": "Generic global rigidity in three dimensions", "category_id": 2, "category": {"id": 2, "name": "combinatorics", "display_name": "Combinatorics", "description": "Counting problems, graph theory, discrete structures.", "slug": "combinatorics", "order_index": 2, "created_at": "2026-07-31T15:26:25.671Z"}, "status": "partially_solved", "problem_number": "AMR-029-0037", "set_id": 13, "statusQualification": "Randomized polynomial recognition is known and stated in the source; the unresolved target is deterministic or combinatorial recognition."}
% FIRST_POSTED: 2026-10-01
% ENTRY_KIND: statement_only
% UnsolvedMath ID 3000037; source catalogue code AMR-029-0037
% !TeX program = lualatex
% Definition and sources only; this file is not an LLM solution attempt.
\documentclass[11pt,a4paper]{article}
\usepackage[margin=25mm]{geometry}
\usepackage{fontspec}
\setmainfont{lmroman10-regular.otf}[BoldFont=lmroman10-bold.otf,
 ItalicFont=lmroman10-italic.otf,BoldItalicFont=lmroman10-bolditalic.otf]
\usepackage{amsmath,amssymb,mathtools}
\usepackage{unicode-math}
\setmathfont{latinmodern-math.otf}
\AtBeginDocument{\let\setminus\smallsetminus}
\usepackage{xurl,hyperref}
\hypersetup{colorlinks=true,urlcolor=blue}
\setcounter{secnumdepth}{0}
\setlength{\parindent}{0pt}
\setlength{\parskip}{5pt}
\setlength{\emergencystretch}{3em}
\begin{document}
\section*{Generic global rigidity in three dimensions}\label{3000037}
{\small UnsolvedMath ID 3000037 · AMR-029-0037 · Combinatorics · AMR Open Problem Lists}
\subsection{Definitions and mathematical statement}
Let $G=(V,E)$ be a finite simple graph. A placement $p:V\to\mathbb R^3$
is generic when all its coordinates are algebraically independent
over $\mathbb Q$. Two placements $p,q$ are equivalent if
\[
                       \|p(u)-p(v)\|=\|q(u)-q(v)\|
                                           \quad(uv\in E),
\]
and congruent if this equality holds for every pair of vertices.
Generic global rigidity means that every placement equivalent to a
generic one is congruent to it. This property depends only on $G$.

The problem is to recognize generically globally rigid graphs in
three-dimensional space by a deterministic polynomial-time algorithm
or an effective combinatorial characterization. The input consists of
vertices and edges, not a framework given with possibly exceptional
coordinates.

The source notes an algebraic characterization giving randomized
polynomial-time recognition. The remaining target is therefore not
bare decidability or the existence of a probabilistic test. Nor is it
enough to test local rigidity, which excludes a continuous flex but can
allow isolated noncongruent equivalent frameworks. Congruence allows
reflections as well as rotations and translations; orientation of the
ambient space is not fixed. Small complete graphs satisfy the definition
without any full-dimensional assumption on their affine spans, so
those cases can be handled directly.
\subsection{Short English statement}
Find deterministic or combinatorial recognition of graphs whose
generic three-dimensional realizations are uniquely determined by their edge lengths.
\subsection{Sources}
\begin{itemize}
\item[S1] ULAM.AI and the UnsolvedMath Contributors, supplied problems.json, record 3000037 (AMR-029-0037), AMR Open Problem Lists. Catalogue identity and source transcription; CC BY 4.0.
\url{https://huggingface.co/datasets/ulamai/UnsolvedMath}.
\item[S2] Egres Open - List of open problems, Open-problem page: Generic global rigidity in three dimensions. Original source indexed by AMR-029-0037.
\url{https://oldlemon.cs.elte.hu/egres/open/List_of_open_problems}.
\item[S3] EGRES Open, Generic global rigidity in three dimensions. Individual source page and explanatory remarks.
\url{https://lemon.cs.elte.hu/egres/open/Generic_global_rigidity_in_three_dimensions}.
\end{itemize}
\end{document}
